\documentclass[aps,prd,preprint,nofootinbib,superscriptaddress,longbibliography]{revtex4-2}

\usepackage{amsmath,amssymb,mathtools,graphicx,bm,booktabs}
\usepackage[colorlinks=true,allcolors=blue]{hyperref}

\newcommand{\phitwo}{\phi^4}
\newcommand{\std}{\mathrm{std}}
\newcommand{\ESS}{\mathrm{ESS}}
\newcommand{\maybegraphic}[3][0.9\linewidth]{%
  \IfFileExists{#2}{\includegraphics[width=#1]{#2}}{%
    \fbox{\parbox[c][2.1in][c]{0.85\linewidth}{\centering \textbf{Figure placeholder}\\[0.5em]#3}}%
  }%
}

\begin{document}

\title{Wilsonian Transport Maps for Two-Dimensional Lattice $\phi^4$ Theory}

\author{Kostas Orginos}

\begin{abstract}
We formulate a normalizing flow for two-dimensional lattice $\phi^4$ theory in
which each renormalization-group blocking step is treated explicitly as a split
into one coarse scalar field and three local fluctuation modes per block, and
only the fluctuation modes are trivialized while the coarse field is held fixed.
This gives the flow a direct Wilsonian interpretation: at each level the model
learns the conditional measure of local fluctuations in the background of the
coarse field. We augment this hierarchy with learned multivariate Gaussian
priors for the fluctuation modes at every RG level and for the terminal
$2\times2$ block. The resulting map is exactly invertible and supports exact
likelihoods and exact reweighting. We further introduce level-resolved
diagnostics based on HMC target samples, including module knockouts and a
conditional decomposition of the model log probability across RG levels. At the
canonical point $(m^2,\lambda)=(-0.4,2.4)$, the best current $16^2$ model
reaches $\std(\Delta S)=0.7869$ and $\ESS=0.5697$, while a direct $32^2$
extension reaches $\std(\Delta S)=1.4740$ and $\ESS=0.2213$ at the end of the
$512$-batch stage. The intensive mismatch $\std(\Delta S)/L$ is nearly
unchanged between $16^2$ and $32^2$, indicating that the residual error behaves
predominantly as an extensive sum of local mismatches rather than as a
catastrophic failure of the transport map. This suggests a useful middle ground
between exact flow-based reweighting and a more Wilsonian program based on
approximate effective actions and universality.
\end{abstract}

\maketitle

\section{Introduction}

Normalizing flows are now a standard tool in lattice field theory for exact
reweighting, accelerated sampling, and transport-map preconditioning
\cite{Dinh2017RealNVP,Albergo2019FlowMCMC,Kanwar2020Equivariant,Albergo2021FlowScaling}.
Much of the existing work constructs invertible maps directly on the full fine
lattice, with inductive bias supplied by symmetries, convolutional structure, or
gauge equivariance. In this work we pursue a different organizing principle: the
transport map should follow a Wilsonian decomposition as directly as possible.

At each $2\times2$ block, the field is decomposed into one coarse mode and
three local fluctuation modes. The coarse field is passed unchanged to the next
level, while only the local fluctuations are trivialized. This makes the model a
learned inverse RG step rather than a generic full-field flow. It also leads to
a clean separation between three conceptually distinct objects:
\begin{enumerate}
\item the deterministic blocking map,
\item the conditional fluctuation trivializer,
\item the effective measure induced on the coarse field.
\end{enumerate}

The present paper focuses on the first two. We show that a Wilsonian
fluctuation-only flow, augmented with learned Gaussian priors at each RG level,
is already accurate enough to expose a nontrivial scaling pattern across
volumes. The main point is not that the architecture is universally superior to
all other hierarchical flows. Rather, the interesting observation is that the
residual mismatch appears to be \emph{intensive}: when the linear size doubles
from $L=16$ to $L=32$, the action-difference width $\std(\Delta S)$
approximately doubles, but $\std(\Delta S)/L$ remains nearly constant.

This observation motivates two complementary interpretations. In the short term,
the model is an exactly invertible transport map with useful reweighting
performance and a direct RG meaning. In the longer term, it suggests that the
learned action may already sit near the target in theory space even when
reweighting becomes volume-limited by extensivity. That viewpoint naturally
connects the present architecture to a more Wilsonian program based on effective
actions and universality.

\section{Wilsonian Transport Map}
\label{sec:wilsonian_map}

Consider an $L\times L$ lattice with $L=2^K$. On each $2\times2$ block we write
\begin{equation}
v=(v_{00},v_{01},v_{10},v_{11})\in\mathbb{R}^4.
\end{equation}
For the default \texttt{average} blocking rule used in our best-performing
models, the code applies the orthonormal Haar transform
\begin{equation}
(c,\eta_1,\eta_2,\eta_3)=Hv,
\end{equation}
with
\begin{equation}
H=
\begin{pmatrix}
\tfrac12 & \tfrac12 & \tfrac12 & \tfrac12 \\
\tfrac12 & -\tfrac12 & \tfrac12 & -\tfrac12 \\
\tfrac12 & \tfrac12 & -\tfrac12 & -\tfrac12 \\
\tfrac12 & -\tfrac12 & -\tfrac12 & \tfrac12
\end{pmatrix}.
\end{equation}
Thus each level-$k$ field is rewritten as a field on the coarse lattice,
\begin{equation}
y_b^{(k)} = \left(c_b^{(k+1)},\eta_b^{(k)}\right),
\qquad
c_b^{(k+1)}\in\mathbb{R},
\qquad
\eta_b^{(k)}\in\mathbb{R}^3.
\end{equation}

The defining feature of the present construction is that, at each nonterminal
level, the coarse field is held fixed while the fluctuation field is mapped to a
simple latent measure. If $q_\theta$ denotes the learned model, the hierarchy
factorizes as
\begin{equation}
q_\theta(\phi^{(0)})=
q_\theta(\phi^{(K-1)})
\prod_{k=0}^{K-2}
q_\theta\!\left(\eta^{(k)}\,\middle|\,\phi^{(k+1)}\right).
\label{eq:hier_factorization}
\end{equation}
Equation~\eqref{eq:hier_factorization} is the sense in which the model is
Wilsonian: the level map learns the conditional measure of the local
fluctuations, while the coarse field carries the effective action to the next
scale.

\begin{figure}[t]
  \centering
  \maybegraphic{figures/wilsonian_hierarchy.pdf}{Schematic of the $2\times2$
  Wilsonian split into one coarse mode and three local fluctuation modes,
  followed by fluctuation trivialization on the coarse lattice.}
  \caption{Schematic view of the Wilsonian transport map. Each blocking step
  decomposes a $2\times2$ block into a coarse scalar and three local
  fluctuations. The coarse field is propagated recursively, while only the
  fluctuations are trivialized.}
  \label{fig:hierarchy}
\end{figure}

\section{Model}
\label{sec:model}

\subsection{Local fluctuation update}

At each nonterminal level, the fluctuation field lives on the coarse lattice and
has three components per site. The update is performed by red/black sweeps on
that coarse lattice. For one red sweep,
\begin{itemize}
\item all coarse scalars are visible and fixed,
\item black fluctuations are visible and frozen,
\item red fluctuations are updated.
\end{itemize}
The black sweep is analogous.

At a coarse-lattice site $b$, the local invertible map is a lower-triangular
affine transformation of the three-component fluctuation vector,
\begin{equation}
\eta_b' = A_b(h)\,\eta_b + t_b(h),
\label{eq:eta_affine}
\end{equation}
where $A_b(h)$ is lower triangular with positive diagonal. The feature $h$ is
produced by a local patch MLP that takes as input the coarse-field patch and the
masked neighboring fluctuation patch. The full map remains triangular across the
checkerboard split, so its Jacobian determinant is exact and cheap.

\subsection{Learned Gaussian priors}

The main improvement over the earlier coarse-eta branch is the addition of a
learned Gaussian whitening layer at each RG level. Before the nonlinear
fluctuation flow, the latent variable
$z_{\eta,b}\sim\mathcal{N}(0,I_3)$ is mapped to
\begin{equation}
\tilde z_{\eta,b}=L^{(\ell)} z_{\eta,b},
\end{equation}
where $L^{(\ell)}$ is a level-dependent lower-triangular matrix with positive
diagonal. In the current best branch this Gaussian is shared across sites within
the same RG level (\texttt{eta\_gaussian = level}).

At the terminal $2\times2$ block, we use a terminal RealNVP together with a
learned zero-mean $4\times4$ Gaussian,
\begin{equation}
z_{\mathrm{term}}
\xmapsto{L_{\mathrm{term}}}
\tilde z_{\mathrm{term}}
\xmapsto{\mathrm{RealNVP}}
x_{\mathrm{term}}.
\end{equation}

All pieces are exactly invertible. If $F_\theta$ denotes the inverse map and
$z=F_\theta(\phi)$, then
\begin{equation}
\log q_\theta(\phi)
=
-\frac12 \|z\|^2
-\frac{N}{2}\log(2\pi)
+ \log\left|\det\frac{\partial F_\theta}{\partial\phi}\right|.
\end{equation}

\subsection{Level-dependent capacity}

The current implementation also allows different nonterminal levels to have
different widths, checkerboard cycle counts, and local patch radii. This matters
numerically because the finest level dominates the aggregate mismatch, whereas
the coarse end of the hierarchy is harder \emph{per degree of freedom}. The
best current $16^2$ model uses
\begin{equation}
\texttt{width\_levels}=[96,64,96],\quad
\texttt{n\_cycles\_levels}=[3,2,3],\quad
\texttt{radius\_levels}=[1,1,2].
\end{equation}

\section{Diagnostics}
\label{sec:diagnostics}

The usual validation quantity is
\begin{equation}
\Delta S(\phi)=S(\phi)+\log q_\theta(\phi),
\end{equation}
up to the sign convention of the training script. The corresponding reweighting
factors are proportional to $e^{-\Delta S}$. In practice we monitor:
\begin{itemize}
\item the centered width $\std(\Delta S-\langle\Delta S\rangle)$,
\item the width of the normalized reweighting factor $w$,
\item an effective-sample-size proxy.
\end{itemize}

To understand \emph{where} the residual error comes from, we also introduced two
level-resolved diagnostics on HMC target samples:
\begin{enumerate}
\item \textbf{module knockouts}: replace selected modules by the identity and
measure the induced degradation in $\std(\Delta S)$,
\item \textbf{conditional decomposition}: measure the per-level conditional
negative log probability,
\(
-\mathbb{E}\bigl[\log q_\theta(\eta^{(k)}\mid \phi^{(k+1)})\bigr]
\),
and the terminal contribution separately.
\end{enumerate}

For a representative $16^2$ learned-Gaussian checkpoint, the knockout ranking
showed that the finest-level nonlinear fluctuation flow is the single most
important module, followed by the terminal flow. The conditional decomposition
showed that the cost per variable increases toward the coarse end of the
hierarchy even though the inverse map already whitens the local latent spaces
well. These diagnostics are important because they tell us that ``where to add
capacity'' is a level-dependent question, not a uniform one.

\section{Results}
\label{sec:results}

All results reported below use the canonical point
\begin{equation}
m^2=-0.4,\qquad \lambda=2.4,
\end{equation}
with explicit $Z_2$ symmetry enforced in the flow.

\subsection{$16^2$}

The best current $16^2$ model uses learned levelwise $3\times3$ Gaussians,
a learned terminal $4\times4$ Gaussian, and the simpler terminal RealNVP
(\texttt{terminal\_n\_layers}=2, \texttt{terminal\_width}=64). After the batch
ramp through $1024$, it reaches
\begin{equation}
\std(\Delta S)=0.7869,\qquad
\std(w)=0.8691,\qquad
\ESS=0.5697.
\end{equation}

\subsection{$32^2$}

For $32^2$ we used a direct extension with one additional fine level and the
per-level schedule
\begin{equation}
\texttt{width\_levels}=[128,96,64,96],\quad
\texttt{n\_cycles\_levels}=[3,3,2,3],\quad
\texttt{radius\_levels}=[1,1,1,2].
\end{equation}
At the end of the $512$-batch stage, the ongoing run has reached
\begin{equation}
\std(\Delta S)=1.4740,\qquad
\std(w)=1.8756,\qquad
\ESS=0.2213.
\end{equation}

The raw mismatch is larger at larger volume, but the more informative intensive
quantity is
\begin{equation}
\frac{\std(\Delta S)}{L}.
\end{equation}
For the current best $16^2$ and current $32^2$ results,
\begin{equation}
\frac{0.7869}{16}=0.0492,\qquad
\frac{1.4740}{32}=0.0461.
\end{equation}
The near constancy of this quantity is the main empirical observation of the
present draft.

\begin{table}[t]
\centering
\begin{tabular}{lcccc}
\toprule
Volume & Stage & $\std(\Delta S)$ & $\std(w)$ & $\ESS$ \\
\midrule
$16^2$ & final ($1024$ batch) & $0.7869$ & $0.8691$ & $0.5697$ \\
$32^2$ & ramp\_512 & $1.4740$ & $1.8756$ & $0.2213$ \\
\bottomrule
\end{tabular}
\caption{Current best $16^2$ result and the current intermediate $32^2$
result. The $32^2$ run is still ongoing, so the second line should be regarded
as preliminary.}
\label{tab:main_results}
\end{table}

\begin{figure}[t]
  \centering
  \maybegraphic{figures/std_deltaS_scaling.pdf}{Planned figure: volume scaling
  of $\std(\Delta S)$ and of the intensive ratio $\std(\Delta S)/L$ for the
  canonical sequence of runs.}
  \caption{Suggested main scaling figure for the paper. The important visual is
  not only that $\std(\Delta S)$ grows with volume, but that the growth is
  close to the expected extensive scaling of a local mismatch.}
  \label{fig:scaling}
\end{figure}

\section{Discussion}
\label{sec:discussion}

The present results suggest a useful reframing of the large-volume problem. A
poor reweighting factor at large volume does not automatically imply that the
learned action is far from the target action in coupling space. If
\begin{equation}
\Delta S(\phi)=\sum_x \delta s_x(\phi)
\end{equation}
is a sum of small local mismatches, then one expects
\begin{equation}
\mathrm{Var}(\Delta S)\propto V,
\qquad
\std(\Delta S)\propto \sqrt{V},
\end{equation}
so exact reweighting becomes more difficult with volume even when the mismatch
remains small in an intensive sense.

This is precisely the pattern suggested by the current $16^2\to32^2$ data. From
the point of view of exact reweighting, the larger-volume run is naturally
harder. From the point of view of theory space, however, the learned action may
already be quite close to the target action. This is why the Wilsonian
formulation matters: it encourages us to think not only in terms of exact
sampling of one microscopic action, but also in terms of learned effective
actions and RG trajectories.

\section{Scope of the Present Paper}
\label{sec:scope}

This draft is intentionally narrow. The novelty claim is \emph{not} that the
present architecture is a universally dominant normalizing flow. The stronger
and more defensible claim is the following:
\begin{quote}
An exactly invertible Wilsonian transport map that trivializes only block
fluctuations, augmented with learned Gaussian priors and level-resolved
diagnostics, produces a small intensive action mismatch whose extensive growth
with volume is predictable and controlled.
\end{quote}

That claim supports a focused near-term paper built around:
\begin{enumerate}
\item the Wilsonian transport-map formulation,
\item the learned local Gaussian priors,
\item the HMC-based level diagnostics,
\item the intensive-scaling interpretation of the residual mismatch.
\end{enumerate}

\section{Outlook}
\label{sec:outlook}

Three immediate continuations suggest themselves.

First, the exact-flow program can be pushed further with better level-dependent
capacity allocation and stage-wise training schedules. The current evidence
already indicates that optimization details matter at least as much as adding
uniform expressivity everywhere.

Second, the present hierarchy naturally suggests a variational RG program. Since
the model explicitly separates blocking from fluctuation trivialization, it
provides a natural starting point for fitting approximate effective actions on
the coarse lattices and tracing RG trajectories in a truncated coupling basis.

Third, the learned map can be used algorithmically as a transport
preconditioner for HMC, for example by performing HMC in latent variables
\cite{Hoffman2019NeuTra} or by embedding the hierarchy into a multilevel update
scheme.

The present paper does not attempt to solve these larger problems. Its purpose
is to establish that Wilsonian transport maps are a useful language for
hierarchical flows in lattice field theory, and that they already reveal a
meaningful intensive-versus-extensive distinction in the residual model error.

\bibliographystyle{apsrev4-2}
\bibliography{refs}

\end{document}
